\documentclass[a4paper,11pt]{article}

% Packages
\usepackage[utf8]{inputenc}
\usepackage[T1]{fontenc}
\usepackage[french]{babel}
\usepackage{amsmath, amssymb, amsthm}
\usepackage{geometry}
\usepackage{xcolor}
\usepackage{tcolorbox}
\usepackage{enumitem}
\usepackage{fancyhdr}

% Configuration de la page
\geometry{margin=2cm}
\pagestyle{fancy}
\fancyhf{}
\fancyhead[L]{Analyse 1}
\fancyhead[R]{Résumé de cours}
\fancyfoot[C]{\thepage}

% Environnements personnalisés
\tcbuselibrary{theorems}

\newtcolorbox{defbox}[1]{
    colback=blue!5!white,
    colframe=blue!75!black,
    fonttitle=\bfseries,
    title=#1
}

\newtcolorbox{propbox}[1]{
    colback=green!5!white,
    colframe=green!75!black,
    fonttitle=\bfseries,
    title=#1
}

\newtcolorbox{theorembox}[1]{
    colback=red!5!white,
    colframe=red!75!black,
    fonttitle=\bfseries,
    title=#1
}

\newtcolorbox{methodbox}[1]{
    colback=orange!5!white,
    colframe=orange!75!black,
    fonttitle=\bfseries,
    title=#1
}

% Titre
\title{\textbf{Analyse 1} \\ Résumé de cours}
\author{}
\date{}

\begin{document}

\maketitle
\tableofcontents
\newpage

\section{Les nombres réels : Généralités et propriétés}

\begin{defbox}{Définition : Majorant minorant borné}
Soit $A$ une partie non-vide de $\mathbb{R}$ :
\begin{itemize}
    \item $M \in \mathbb{R}$ est \textbf{majorant} de $A$ si $\forall a \in A, \; a \leq M$
    \item $m \in \mathbb{R}$ est \textbf{minorant} de $A$ si $\forall a \in A, \; a \geq m$
    \item $A$ est \textbf{bornée} si elle est à la fois majorée et minorée
\end{itemize}
\end{defbox}

\begin{defbox}{Définition : Maximum et minimum}
\begin{itemize}
    \item $M$ est le \textbf{maximum} de $A$ si $M$ est un majorant de $A$ et $M \in A$
    \item $m$ est le \textbf{minimum} de $A$ si $m$ est un minorant de $A$ et $m \in A$
\end{itemize}
\end{defbox}

\begin{defbox}{Définition : Supremum et infimum}
\begin{itemize}
    \item Si $A$ majorée : le \textbf{supremum} est le plus petit des majorants (noté sup $A$)
    \item Si $A$ minorée : l'\textbf{infimum} est le plus grand des minorants (noté inf $A$)
\end{itemize}
\end{defbox}

\section{Rappel sur les fonctions}

\subsection{Fonctions croissante, décroissante, monotone}

\begin{defbox}{Définition}
Soit $f$ une fonction et $D \subseteq \mathbb{R}$ :
\begin{itemize}
    \item $f$ \textbf{croissante} sur $D$ : $\forall x, y \in D, \; x < y \Rightarrow f(x) \leq f(y)$
    \item $f$ \textbf{décroissante} sur $D$ : $\forall x, y \in D, \; x < y \Rightarrow f(x) \geq f(y)$
    \item $f$ \textbf{monotone} : si croissante ou décroissante
\end{itemize}
\end{defbox}

\subsection{Fonctions paires et impaires}

\begin{defbox}{Définition : Domaine symétrique}
$D$ symétrique par rapport à 0 : $\forall x \in D, \; -x \in D$

Sur un tel domaine :
\begin{itemize}
    \item $f$ \textbf{paire} : $\forall x \in D, \; f(-x) = f(x)$
    \item $f$ \textbf{impaire} : $\forall x \in D, \; f(-x) = -f(x)$
\end{itemize}
\end{defbox}

\subsection{Fonctions périodiques}

\begin{defbox}{Définition : Fonction $T$-périodique}
$f$ est $T$-périodique ($T > 0$) si : $\forall x \in D, \; f(x + T) = f(x)$
\end{defbox}

\newpage
\section{Fonctions polynomiales et rationnelles}

\subsection{Fonction polynomiale d'ordre $n$}

\begin{defbox}{Définition}
$P(x) = a_0 + a_1x + \ldots + a_nx^n$ avec $a_n \neq 0$
\end{defbox}

\subsection{Fonction rationnelle}

\begin{defbox}{Définition}
$F(x) = \frac{P(x)}{Q(x)}$ où $P$ et $Q$ sont des fonctions polynomiales
\end{defbox}

\section{Fonctions trigonométriques}

\subsection{Cercle trigonométrique}

\begin{propbox}{Propriétés importantes}
\begin{itemize}
    \item $M(a) \in \mathbb{R}$ correspond au point $M$ du cercle unité
    \item Valeurs usuelles : 
    \begin{itemize}
        \item $\sin(0) = 0$, $\cos(0) = 1$
        \item $\sin(\frac{\pi}{2}) = 1$, $\cos(\frac{\pi}{2}) = 0$
        \item $\sin(\pi) = 0$, $\cos(\pi) = -1$
    \end{itemize}
\end{itemize}
\end{propbox}

\subsection{Formules d'addition}

\begin{propbox}{Formules essentielles}
\begin{align*}
\cos(a + b) &= \cos a \cos b - \sin a \sin b \\
\sin(a + b) &= \sin a \cos b + \cos a \sin b \\
\cos(a - b) &= \cos a \cos b + \sin a \sin b \\
\sin(a - b) &= \sin a \cos b - \cos a \sin b
\end{align*}
\end{propbox}

\newpage
\section{Limites et continuité}

\subsection{Limites : rappels}

\begin{defbox}{Définition : Taux de variation}
$T_f(a,b) = \frac{f(b) - f(a)}{b - a}$
\end{defbox}

\subsection{Théorème de l'encadrement (Théorème des gendarmes)}

\begin{theorembox}{Théorème}
Si $g(x) \leq f(x) \leq h(x)$ et $\lim_{x \to a} g(x) = \lim_{x \to a} h(x) = \ell$

Alors $\lim_{x \to a} f(x) = \ell$
\end{theorembox}

\subsection{Méthodes classiques de calcul de limites}

\subsubsection{Limite d'une fonction polynomiale à l'infini}

\begin{methodbox}{Méthode}
Pour $P(x) = a_0 + a_1x + \ldots + a_nx^n$ avec $a_n \neq 0$ :

\[\lim_{x \to \pm\infty} P(x) = \lim_{x \to \pm\infty} a_nx^n\]

\textbf{Cas $n$ pair :}
\begin{itemize}
    \item Si $a_n > 0$ : $\lim_{x \to \pm\infty} P(x) = +\infty$
    \item Si $a_n < 0$ : $\lim_{x \to \pm\infty} P(x) = -\infty$
\end{itemize}

\textbf{Cas $n$ impair :}
\begin{itemize}
    \item Si $a_n > 0$ : $\lim_{x \to +\infty} P(x) = +\infty$ et $\lim_{x \to -\infty} P(x) = -\infty$
    \item Si $a_n < 0$ : $\lim_{x \to +\infty} P(x) = -\infty$ et $\lim_{x \to -\infty} P(x) = +\infty$
\end{itemize}
\end{methodbox}

\subsubsection{Limite d'une fonction rationnelle à l'infini}

\begin{methodbox}{Méthode}
Pour $F(x) = \frac{P(x)}{Q(x)}$ avec $P$ de degré $n$ et $Q$ de degré $m$ :

\begin{itemize}
    \item Si $n = m$ : $\lim_{x \to \pm\infty} F(x) = \frac{a_n}{b_m}$
    \item Si $m > n$ : $\lim_{x \to \pm\infty} F(x) = 0$
    \item Si $n > m$ : la limite est $\pm\infty$ selon le signe de $\frac{a_n}{b_m}$ et la parité de $n - m$
\end{itemize}
\end{methodbox}

\subsubsection{Limite en un pôle (quand $Q(a) = 0$)}

\begin{methodbox}{Méthode}
Pour $F(x) = \frac{P(x)}{Q(x)}$ avec $Q(a) = 0$ :

On factorise : $Q(x) = (x - a)^p S(x)$ avec $S(a) \neq 0$

De même pour $P$ si $P(a) = 0$

Alors : $F(x) = (x - a)^m \frac{R(x)}{S(x)}$ avec $m \in \mathbb{Z}$

\textbf{Cas :}
\begin{itemize}
    \item Si $m = 0$ : $\lim_{x \to a} F(x) = \frac{R(a)}{S(a)}$
    \item Si $m < 0$ : étudier séparément $m$ pair et impair
\end{itemize}
\end{methodbox}

\subsubsection{Forme indéterminée avec racines}

\begin{methodbox}{Méthode : Quantité conjuguée}
Pour calculer $\lim_{x \to +\infty} (\sqrt{x + 3} - \sqrt{x + 1})$ :

\textbf{Astuce :} multiplier et diviser par la quantité conjuguée

\[f(x) = \frac{(\sqrt{x + 3} - \sqrt{x + 1})(\sqrt{x + 3} + \sqrt{x + 1})}{\sqrt{x + 3} + \sqrt{x + 1}} = \frac{2}{\sqrt{x + 3} + \sqrt{x + 1}}\]

Donc $\lim_{x \to +\infty} f(x) = 0$
\end{methodbox}

\subsection{Théorème de L'Hôpital}

\begin{theorembox}{Théorème de L'Hôpital}
Soit $f$ et $g$ deux fonctions dérivables au voisinage de $a$ (sauf peut-être en $a$).

Si on a une forme indéterminée $\frac{0}{0}$ ou $\frac{\infty}{\infty}$ et que $\lim_{x \to a} \frac{f'(x)}{g'(x)}$ existe, alors :

\[\lim_{x \to a} \frac{f(x)}{g(x)} = \lim_{x \to a} \frac{f'(x)}{g'(x)}\]

\textbf{Attention :} Ce théorème s'applique aussi pour $a = \pm\infty$
\end{theorembox}

\begin{methodbox}{Application pratique}
\textbf{Exemple :} Calculer $\lim_{x \to 0} \frac{\sin x}{x}$

On a la forme indéterminée $\frac{0}{0}$

En appliquant L'Hôpital : $\lim_{x \to 0} \frac{\sin x}{x} = \lim_{x \to 0} \frac{\cos x}{1} = 1$
\end{methodbox}

\newpage
\section{Continuité}

\subsection{Prolongement par continuité}

\begin{defbox}{Définition}
Soit $f : I \backslash \{x_0\} \to \mathbb{R}$ et $\ell = \lim_{x \to x_0} f(x)$

On définit $\tilde{f}(x) = \begin{cases} f(x) & \text{si } x \neq x_0 \\ \ell & \text{si } x = x_0 \end{cases}$

$\tilde{f}$ s'appelle le \textbf{prolongement par continuité} de $f$ en $x_0$
\end{defbox}

\subsection{Théorème des valeurs intermédiaires (TVI)}

\begin{theorembox}{Théorème des valeurs intermédiaires}
Soit $f$ continue sur $[a, b]$. Alors $f$ atteint toutes les valeurs comprises entre $f(a)$ et $f(b)$ :

\[\forall y \in [\min(f(a), f(b)), \max(f(a), f(b))], \; \exists x \in [a, b] : y = f(x)\]
\end{theorembox}

\subsection{Théorème des bornes atteintes (TBA - Admis)}

\begin{theorembox}{Théorème des bornes atteintes}
Soit $f : [a, b] \to \mathbb{R}$ continue sur $[a, b]$. 

Alors $f$ est bornée sur $[a, b]$ et atteint ses bornes :

\[\exists \alpha, \beta \in [a, b] : f(\alpha) = \min_{x \in [a,b]} f(x) \text{ et } f(\beta) = \max_{x \in [a,b]} f(x)\]
\end{theorembox}

\newpage
\section{Fonctions monotones et bijections}

\subsection{Bijection}

\begin{defbox}{Définition : Bijection}
$f : D \to \mathbb{R}$ est une \textbf{bijection} de $D$ sur $D'$ si :
\begin{enumerate}
    \item $f(D) \subset D'$ (image incluse dans $D'$)
    \item $\forall y \in D', \; \exists! x \in D : y = f(x)$ (tout élément de $D'$ a un unique antécédent)
\end{enumerate}
\end{defbox}

\subsection{Bijection réciproque}

\begin{defbox}{Définition : Bijection réciproque}
Si $f : D \to D'$ est une bijection, la \textbf{bijection réciproque} $g : D' \to D$ vérifie :

\[\forall y \in D', \; x = f^{-1}(y) \Leftrightarrow \begin{cases} x \in D \\ y = f(x) \end{cases}\]

On note $g = f^{-1}$
\end{defbox}

\subsection{Théorème de la bijection}

\begin{theorembox}{Théorème}
Soit $f : I \to \mathbb{R}$ continue et strictement monotone sur $I = (a, b)$. Alors :
\begin{enumerate}
    \item $f$ établit une bijection de $I$ dans $J = f(I)$
    \item La bijection réciproque $f^{-1}$ est continue et strictement monotone sur $J$, avec le même sens de variation que $f$
\end{enumerate}
\end{theorembox}

\subsection{Exemples de fonctions réciproques}

\subsubsection{Fonction carré et racine carrée}

\begin{propbox}{Propriété}
$f : [0, +\infty[ \to \mathbb{R}$ définie par $f(x) = x^2$ est une bijection de $[0, +\infty[$ sur $[0, +\infty[$

La bijection réciproque est $f^{-1} : x \mapsto \sqrt{x}$
\end{propbox}

\subsubsection{Fonction cube et racine cubique}

\begin{propbox}{Propriété}
$f : \mathbb{R} \to \mathbb{R}$ définie par $f(x) = x^3$ est une bijection de $\mathbb{R}$ sur $\mathbb{R}$

La bijection réciproque est $f^{-1}(x) = \sqrt[3]{x}$ ou $x^{1/3}$
\end{propbox}

\subsubsection{Racine $n$-ième}

\begin{propbox}{Propriété : $n$ pair strictement positif}
$f : [0, +\infty[ \to \mathbb{R}$ définie par $f(x) = x^n$ est une bijection sur $[0, +\infty[$

La bijection réciproque est $f^{-1}(x) = \sqrt[n]{x} = x^{1/n}$

Pour $n = 2$ : $\sqrt[2]{x} = \sqrt{x}$
\end{propbox}

\begin{propbox}{Propriété : $n$ impair positif}
$f : \mathbb{R} \to \mathbb{R}$ définie par $f(x) = x^n$ est une bijection sur $\mathbb{R}$

La bijection réciproque est $f^{-1}(x) = \sqrt[n]{x} = x^{1/n}$
\end{propbox}

\newpage
\section{Fonctions circulaires réciproques}

\subsection{Rappels}

\begin{propbox}{Formules importantes}
\begin{itemize}
    \item $\forall x \in [-1, 1], \; \cos(\arccos x) = x$ et $\sin(\arcsin x) = x$
    \item $\forall x \in \mathbb{R}, \; \tan(\arctan x) = x$
\end{itemize}

Et réciproquement :
\begin{itemize}
    \item $\arccos(\cos x) = x \Leftrightarrow x \in [0, \pi]$
    \item $\arcsin(\sin x) = x \Leftrightarrow x \in [-\frac{\pi}{2}, \frac{\pi}{2}]$
    \item $\arctan(\tan x) = x \Leftrightarrow x \in ]-\frac{\pi}{2}, \frac{\pi}{2}[$
\end{itemize}
\end{propbox}

\vspace{1cm}
\begin{center}
\textit{--- Fin du résumé ---}
\end{center}

\end{document}
